\begin{figure}[htbp]
\centering
\begin{tikzpicture}[scale=.92,>=Stealth,font=\small]
  \coordinate (O) at (0,0);
  \coordinate (I) at (3.5,-.65);
  \coordinate (J) at (-3,-.9);
  \coordinate (K) at (0,2.65);
  \fill[blue!9] (O)--(-2.35,-.705)--(-2.35,1.35)--(0,2.055)--cycle;
  \fill[orange!14] (O)--(2.65,-.492)--(2.65,1.563)--(0,2.055)--cycle;
  \draw[->,semithick] (O)--(I) node[right,align=left] {$[\omega_I]$\\$\in\R_{>0}h$};
  \draw[->,semithick] (O)--(J) node[left] {$[\omega_J]$};
  \draw[->,semithick] (O)--(K) node[above] {$[\omega_K]\in\R_{>0}(v-a\lambda e)$};
  \node[align=center] at (-1.27,.78) {period plane\\for $I$};
  \node[align=center] at (1.42,.90) {period plane\\for $J$};
  \fill (O) circle (1.2pt);
\end{tikzpicture}
\caption{The period planes before and after rotation, inside the
positive three-space spanned by the hyperk\"ahler forms.
The class $e$ is perpendicular to $[\omega_I]$ and $[\omega_K]$,
so it becomes a $(1,1)$ class for $J$; its pairing with
$[\omega_J]$ fixes the positive sign.}
\label{fig:rotation}
\end{figure}
